% Folder 156-9, batch 06 — archivists' pages 101–120.
% Pass: Opus 5.5 (claude-opus-5-5), 2026-10-03 — first pass, unchecked against the pages by a human.
%
% Pages 109, 111 and 113 are typed versos (notes of another text) carrying no
% mathematics of this run; they are skipped. Grothendieck's own pagination runs
% 98–115 on the transcribed pages (his 108 is overwritten as 109 on p. 114).
\documentclass[11pt,a4paper]{article}
\input{../preamble/grothendieck}

\folder{156-9}
\batch{6}
\pages{101}{120}
\foldertitle{[Chapitre] IX et IX bis. [Ateliers] : notes manuscrites (05-15/07/1986).}
\dating{1986}
\watermark{Édition de démonstration}

\begin{document}

\page{101}
\note{p. 98 de l'auteur. La page s'ouvre au milieu d'une discussion des correspondances entre spatios $\Sigma \to \Sigma'$, commencée avant ce lot ; la relation d'ordre qui suit porte sur $\mathrm{Corr}(\Sigma,\Sigma')$. Notations du lot : $\mathbin{|\circ|}$ pour la relation de disjonction, $\complement$ pour le complémentaire, $\mathfrak{P}(E)$ pour l'ensemble des parties.}

\[
  f \le g \;\overset{\text{déf}}{\iff}\; \forall x \in \Sigma,\ f(x) \le g(x) .
\]

Je dis que pour cette relation d'ordre dans $\mathrm{Corr}(\Sigma,\Sigma')$, les Sup quelconques (donc aussi les Inf quelconques) existent, et sont donnés par
\[
  \bigl(\operatorname{Sup}_i f_i\bigr)(x) = \operatorname{Sup}_i f_i(x) \qquad \forall x \in \Sigma .
\]
En effet, l'associativité des Sup montre que l'application des Sup \uncertain{ainsi} \uncertain{définie} commute bien aux Sup quelc., donc est une corr. de spatios.

Mais la structure $\complement$, ou la structure de disposition, n'est pas claire. Dans le cas où $E$ est une disposante triviale, \struck{\ill{}} et où les corr. $E \to \Sigma'$ sont \uncertain{donc} les applications quelconques $E \to \Sigma'$, donc
\[
  \mathrm{Corr}(\Sigma_E, \Sigma') \simeq \Sigma'^{E},
\]
des structures de disjonction et de $\complement$ \uncertain{se définissent} « terme par terme », \uncertain{prenant} \ill{}
\[
  (\complement f)(x) = \complement f(x) \qquad \forall x \in E ,
\]

\page{102}
\note{p. 99 de l'auteur.}
\[
  f \mathbin{|\circ|} g \iff \forall x \in E,\ f(x) \mathbin{|\circ|} g(x) ,
\]
mais si on prend des \uncertain{arguments} dans $\mathfrak{P}(E)$ et non dans $E$, on n'a \uncertain{plus} en général
\[
  (\complement f)(A) \ne \complement\bigl(f(A)\bigr) \qquad \bigl(A \in \Sigma_E = \mathfrak{P}(E)\bigr)
\]
et
\[
  f \mathbin{|\circ|} g \not\Longrightarrow f(A) \mathbin{|\circ|} g(A) .
\]

On voit donc que la question de définir $\complement$ ou $\mathbin{|\circ|}$ sur $\mathrm{Corr}(\Sigma,\Sigma')$ en général, semble liée à la question de l'\uncertain{existence} \ill{} le droit, \struck{\ill{}} des éléments minimaux de $\Sigma \smallsetminus \{0_\Sigma\}$, ou d'autres objets \uncertain{arithmétiques} qui en tiendraient lieu, jouant le rôle des « points ».

Une \uncertain{question} \ill{} dans \ill{} \uncertain{l'aristotélicienne}. \ill{} \uncertain{opération} est la suivante. \struck{Dans le} \ill{} ensembliste (disp. triviales), utilisant le graphe $\Gamma$ \add{$(\subset E \times E')$} d'une correspondance, \struck{celle-ci peut}

\page{103}
\note{p. 100 de l'auteur.}

\begin{tikzcd}
  \Gamma \arrow[r, "f'"] \arrow[d, "f"'] & E' \\
  E &
\end{tikzcd}

celle-ci peut se décomposer en
\[
  \Gamma = f' \circ {}^{t}f ,
\]
i.e. « l'inverse » d'une application (\uncertain{correspondant} \uncertain{à} : $A \mapsto f^{-1}(A)$, $\mathfrak{P}(E) \to \mathfrak{P}(\Gamma)$) suivie de l'application $f'$ (\uncertain{donnant} \ill{} : $B \mapsto f'(B)$, $\mathfrak{P}(\Gamma) \to \mathfrak{P}(E')$).

Une corr. de la forme ${}^{t}f$ \struck{\ill{}} (« comorphisme ») peut d'ailleurs se \uncertain{caractériser} en termes de corr. de spatios \uncertain{algébriquement} \struck{et \ill{}}, par la propriété de commutation à $\complement$ ; et de même, une correspondance de la forme $f'$ (« morphisme ») se caractérise par la propriété que sa transposée ${}^{t}f'$ commute à $\complement$.

Y a-t-il une décomposition similaire dans le cas général des \struck{corr. de} \uncertain{atelier} \uncertain{dans} \ill{} disposantes, ou en qui revient au même, entre \uncertain{spatios} ?

\page{104}
\note{p. 101 de l'auteur.}
Ces questions ne \uncertain{paraissent} pas, à \uncertain{première} \uncertain{vue}, être de conséquence directe pour la formalisation ensembliste des ateliers. Il n'en va plus de même des suivantes.

\textbf{Analogie ensembliste.} Soient $E$, $E'$ des ensembles, et
\[
  \alpha_*\colon \mathfrak{P}(E) \to \mathfrak{P}(E')
\]
une application associée à une corr. $\alpha$, i.e. $\alpha_*$ commutant aux Sup quelc. Soit $\Gamma$ le graphe, \struck{\ill{}} $E \xleftarrow{f} \Gamma \xrightarrow{f'} E'$, donc $\alpha_*(A) = f'\bigl(f^{-1}(A)\bigr)$.
\note{Sur la plupart des $\alpha_*$ des pp. 101--105, un petit trait griffonné surmonte la lettre $\alpha$ ; on ne sait s'il s'agit d'un accent biffé ou d'un signe distinctif, et il n'est pas reproduit.}

\textbf{A)} Conditions équivalentes
\begin{itemize}
  \item[(i)] $\alpha_*$ commute à $\mathbin{|\circ|}$ ;
  \item[(ii)] $\alpha_*$ commute à : $S \wedge T$, $\bigwedge_{i\in I} S_i$, pour $I \ne \emptyset$ ;
  \item[(ii$'$)] $\alpha_*$ commute à : $S \wedge \complement T$ ;
  \item[(iii$'$)] $\alpha_*$ commute à : $S \wedge \complement T$ pour $T \le S$ ;
  \item[(iv)] $f'$ injective --- i.e. $\alpha$ composé d'un « comorphisme », et d'un « morphisme » associé à une application injective (ou \uncertain{encore} : \ill{} \uncertain{définissant} une application $f'$ injective $\mathfrak{P}(\Gamma) \to \mathfrak{P}(E')$.
\end{itemize}

\textbf{Corollaire 1.} Conditions équivalentes
\begin{itemize}
  \item[(i)] $\alpha_*$ commute à $\mathbin{|\circ|}$ et $\alpha_*(1) = 1$ ;
  \item[(ii)] $\alpha_*$ --- \,$S \wedge T$ et $\alpha_*(1) = 1$ ;
  \item[(ii$'$)] $\alpha_*$ --- \,$\bigwedge_{i\in I} S_i$, $I$ fini.
\end{itemize}
\marginal{\uncertain{déduit} \uncertain{de} A et B, \uncertain{cor.} \ill{}}

\page{105}
\note{p. 102 de l'auteur.}
\begin{itemize}
  \item[(ii$''$)] $\alpha_*$ commute aux $\bigwedge$ quelconques ;
  \item[(iii)] $\alpha_*$ commute à : $\complement$, i.e. $f\colon \Gamma \xrightarrow{\sim} E$ ;
  \item[(iv)] $\alpha$ est un comorphisme, i.e. \struck{$\alpha_*$} \add{$\alpha_*$} de la forme $A \mapsto g^{-1}(A)$, où $g \;(= f' f^{-1})\colon E' \to E$.
\end{itemize}
\note{la parenthèse $g = f' f^{-1}$ est lue telle quelle ; le sens de la flèche $E' \to E$ demanderait plutôt $f \circ f'^{-1}$.}

\textbf{B)} Conditions équivalentes
\begin{itemize}
  \item[(i)] $\alpha_*(1_E) = 1_{E'}$ ;
  \item[(ii)] $\beta_*(S) = \emptyset_E \Longrightarrow S = \emptyset_{E'}$ ;
  \item[(iii)] $f'$ surjectif --- i.e. $\alpha$ composé d'un comorphisme, et d'un morphisme surjectif.
\end{itemize}

\textbf{Corollaire 2.} Conditions équivalentes
\begin{itemize}
  \item[(i)] $\alpha_*$ commute à : $\mathbin{|\circ|}$ « fidèlement », i.e. $\forall S, T \in \mathfrak{P}(E)$,
  \[
    S \mathbin{|\circ|} T \iff \alpha(S) \mathbin{|\circ|} \alpha(T) ;
  \]
  \item[(ii)] $\alpha_*$ commute à : $\mathbin{|\circ|}$ (cf.\ A) \ill{} conditions équivalentes : celles-ci), et
  \[
    \alpha_*(S) = \emptyset \Longrightarrow S = \emptyset ;
  \]
  \item[(iii)] $\alpha$ est composé d'un comorphisme défini par une application surjective $g$, et d'un morphisme injectif.
\end{itemize}

\textbf{Corollaire 3.} Conditions équivalentes (cf.\ A) pour cond. équiv.),
\begin{itemize}
  \item[(i)] $\alpha_*$ commute à : $\complement$, et $\alpha_*(S) = 0_{E'} \Rightarrow S = 0_E$ ;
  \item[(ii)] $\alpha_*$ commute à : $\mathbin{|\circ|}$ fidèlement, et $1 \to 1$ ;
  \item[(iii)] $\alpha$ est un comorphisme correspondant à une \add{application} surjective $g\colon E' \to E$.
\end{itemize}

\page{106}
\note{p. 103 de l'auteur.}
La question se pose donc dans quelle mesure ces équivalences \struck{sont valables} gardent un sens et restent valables pour des correspondances entre disposantes quelconques. Faute d'avoir une décomposition \uncertain{canonique} du type $\alpha_* = f'_* f^*$, \struck{il y a} l'\uncertain{interprétation} « naïve » ou « directe » en termes des propriétés de $f$, $f'$, \struck{\ill{}} semble \uncertain{devoir} disparaître. En fait, on a des équivalents algébriques « \uncertain{naturels} » \uncertain{de} \ill{}, qui ont un sens \ill{}.

Mathématiquement, \ill{} \uncertain{les} \uncertain{preuves} \add{(\ill{} que (i) \ill{} (iii) dans A)} des équivalences de (i) \ill{} (iii$'$) dans A), \uncertain{qui} \uncertain{passent} par des hypothèses draconiennes, \uncertain{plus} \uncertain{fortes} \uncertain{que} de l'hypothèse booléenne \ill{} \ill{} des \uncertain{opérations} à \ill{} \ill{}. Peut-être \ill{} \ill{} des faux problèmes, en ce que pour mon travail proprement géométrique, \uncertain{on} n'aura \uncertain{affaire} qu'à des supports très particuliers (constructibles), les supports \ill{} (\uncertain{constructibles}) et on \uncertain{n'aura} à prendre des Sup et des Inf que \uncertain{sous} des conditions idoines de « commensurabilité ». Ce sont \uncertain{les} \uncertain{cas} \ill{}, précisément, les arguments
\marginal{\uncertain{cf.\ VIII} \ill{}}

\page{107}
\note{p. 104 de l'auteur.}
auxquels j'ai fait allusion tout à l'heure, dans A), pour nous dire que la commutation de $\alpha_*$ à $\mathbin{|\circ|}$ implique les autres propriétés (sous les hyp. restrictives de commensurabilité). Dans l'intuition géométrique : « $\alpha_*$ commute à $\mathbin{|\circ|}$ » \uncertain{signifie} \uncertain{qu'elle} \uncertain{est} celle d'un « comorphisme » \uncertain{suivi} d'une application injective. Quand de plus $\alpha_*(1) = 1$, on s'attend donc à avoir commutation à $\complement$ (\ill{} \ill{} le \uncertain{prouver}, moyennant des conditions de constructibilité convenables).

\textbf{NB.} La commutation à $\complement$ \uncertain{implique} \uncertain{déjà} que $1 \to 1$, ainsi que la compatibilité à $\mathbin{|\circ|}$.

\textbf{Question.} Si $\alpha_*$ et $\beta_*$ commutent à $\complement$ (\uncertain{i.e.} \ill{}), \ill{} un \uncertain{iso.} de spatios, \uncertain{inverses} \uncertain{l'un} de l'autre) ?

Notons enfin que l'équivalence de (i) et (ii) dans (B) se vérifie trivialement dans tout spatios.

\page{108}
\note{p. 105 de l'auteur.}
\textbf{Lemme 1.} Soit $b \le c$ dans $\Sigma$. Conditions équivalentes
\begin{itemize}
  \item[(i)] $\exists\, b^* \in \Sigma$ \uncertain{tel que} $b^* \mathbin{|\circ|} b$ et $b \vee b^* = c$ ;
  \item[(ii)] On a $b \vee b^*_0 = c$, où $b^*_0 = \complement b \wedge c$.
\end{itemize}
De plus, si on a $b^*$ comme dans (i), on a $b^* \le b^*_0$, i.e.\ $b^*_0$ est le plus grand des $b^*$ \uncertain{satisfaisant} (i).

\textbf{Lemme 2.} Soit $b \le c$ dans $\Sigma$. Conditions équivalentes (pour $\vee$, $\wedge$, $\complement$),
\begin{itemize}
  \item[(i)] $b$, $c$ engendrent dans $\Sigma$ une sous-algèbre de Boole.
  \item[(ii)] $b$, $c$ engendrent une sous-algèbre de Boole finie, de cardinal $2^n$ avec $n \in \{1,2,3\}$.
\end{itemize}
\add{NB} ($n = 1$ ssi $\{b,c\} \subset \{0,1\}$ ; $n = 2$ ssi $b = 0$ ou $c = 1$ ou $b = c$, \struck{mais $\{b,c\} \not\subset \{0,1\}$ ; $\{b,c\} \ne \{0,1\}$ (\ill{} exclus) $c \ne 0$, $b \ne 1$ et} ; $n = 3$ ssi \struck{\ill{} $b = 0$, $b \ne 1$, et $b = 0$ et $c = 1$}
\[
  0 < b < c < 1 \bigr) .
\]
\begin{itemize}
  \item[(iii)] $(b, c)$ satisfait la condition du lemme 1, et $(c', b')$ aussi, où $c' = \complement c$, $b' = \complement b$.
  \item[(iv)] Posant $\alpha = b$, $\beta = c \wedge \complement b$, $\gamma = \complement c$, \struck{$b$,} $\alpha$, $\beta$, $\gamma$ engendrent une sous-alg. de Boole \ill{} \ill{} de $\Sigma$, \uncertain{satisfaisant} aux relations $\complement(\alpha \vee \beta) = \gamma$, $\complement(\beta \vee \gamma) = \alpha$, \add{$\complement(\gamma \vee \alpha) = \beta$} / i.e.\
  \[
    \complement(\alpha) \wedge (\gamma \vee \alpha) = \gamma, \quad
    \complement(\beta)(\gamma \vee \alpha) = \alpha .
  \]
\end{itemize}
\note{la dernière ligne est lue telle quelle : les lettres et le signe $\wedge$ manquant du second terme ne s'accordent pas avec les relations qu'elle prétend reformuler.}
\marginal{\uncertain{dont} $\alpha, \beta, \gamma$ \uncertain{mut.} disjoints}

\textbf{Cor.} Conditions équivalentes sur $\Sigma$
\begin{itemize}
  \item[(i)] $\forall$ couple $b, c$ avec $b \le c$, les conditions du lemme 1 sont satisfaites. (plus forte)
  \item[(ii)] $\forall$ couple $b, c$ avec $b \le c$, les conditions du lemme 2 sont satisfaites.
  \item[(iii)] $\forall$ couple $\alpha, \beta$ avec $\alpha \mathbin{|\circ|} \beta$, $\alpha$ et $\beta$ engendrent une sous-algèbre de Boole (laquelle sera finie, de card.\ $2^n$, \struck{$1 \le n \in \{1,2,3\}$} \add{ou $\alpha \vee \beta = 1$}, \ill{} ; $n = 1$ ssi $\{\alpha, \beta\} \subset \{0,1\}$, $n \le 2$ ssi $\alpha$ ou $\beta$ est nul \struck{\ill{}}, $n = 3$ ssi $\alpha \ne 0$, $\beta \ne 0$, $\alpha \vee \beta \ne 1$.
  \item[(iiii)] \struck{\ill{} $\alpha \ne 0$,} Condition (iii), mais avec familles finies $(\alpha_i)$, avec $\alpha_i \mathbin{|\circ|} \alpha_j$ \uncertain{pour} $i \ne j$.
\end{itemize}
\marginal{\ill{} \uncertain{lemme} \ill{} (iii) \ill{} \uncertain{plus} \uncertain{cor.\ 2} \ill{}}
\note{la marge gauche porte, en travers, une longue note de quatre ou cinq lignes en partie biffée, dont on ne lit que quelques mots ; elle n'est pas transcrite.}

\page{110}
\note{p. 106 de l'auteur.}
\add{Sous les conditions éq.\ du Cor.\ 1 (iii),}

\textbf{Corollaire 2.} Si $\alpha, \beta \ne 0$, $\alpha \vee \beta \ne 1$, et posant $\gamma = \complement(\alpha \vee \beta) \ne 0$, $\alpha$, $\beta$, $\gamma$ forment une « base » de l'algèbre de Boole engendrée par $\alpha$, $\beta$. Sous les conditions du lemme 2, \add{et supposant $0 < b < c < 1$}, posant
\[
  \alpha = b, \quad \beta = c \wedge \complement b, \quad \gamma = \complement c ,
\]
$\alpha$, $\beta$, $\gamma$ forment une base de l'algèbre de Boole engendrée par $b$, $c$.

\textbf{Question.} \uncertain{Les} conditions envisagées \uncertain{dans} \struck{Lemme 1} Cor.\ 1 sur $\Sigma$ sont-elles \uncertain{toujours} vérifiées ? J'en doute fort ! \uncertain{Le} \uncertain{sont-elles} si $\Sigma$ est fini ? J'en doute également. Si on suppose ces conditions satisfaites, $\Sigma$ est-il nécessairement une algèbre de Boole ? \uncertain{Est-il} \uncertain{du} \uncertain{moins} \uncertain{vrai} l'hypothèse similaire à celle du Cor.\ \add{(iiii)}, \uncertain{pour} les familles $(\alpha_i)_{i \in I}$ d'élts mutuellement disjoints de $I$ (\uncertain{presque} \uncertain{finies}) ?
\marginal{\uncertain{contrexemple} \ill{} $\Sigma = 10$}

Il faudrait des \uncertain{relations} \uncertain{convenables} \struck{\ill{}} \ill{} qui \ill{}, \uncertain{pour} un spatios $\Sigma$ \uncertain{construit} via un ens.\ ordonné $L$, mais \uncertain{où} la relation $x \mathbin{|\circ|} y$ ssi $x < y$ ou $y < x$.

\page{112}
\note{p. 107 de l'auteur.}
\drawing{en haut à gauche, un treillis dessiné en perspective, comme un prisme hexagonal vu de biais : de $0$ partent des flèches vers $a$ (en passant par $a'$), $c$ et $b'$ ; de là vers $B$, $C$, $A$ (avec une étape marquée $A'$ entre $A$ et $1$), puis vers $1$ ; plusieurs flèches intérieures repassées et une étoile au centre. Au-dessous, un hexagone orienté de sommets $A$, $b$, $C$, $a$, $B$, $c$, les flèches allant des sommets bas $b$, $a$, $c$ vers les sommets $A$, $C$, $B$, doublé d'un arc de cercle.}

Points $i$, $i'$, $j$, $k$ (dessinés à côté).
\[
  \text{\struck{$A \in \mathfrak{t}$}} \quad
  E \in \mathfrak{P}^3(I) \;\Bigm|\;
  \left\{
  \begin{array}{l}
    i, i' \in E \text{ ou } i, i' \notin E \text{ ou} \\
    E = \{i\} \text{ ou } E = \{i, j, k\}
  \end{array}
  \right\}
\]
\struck{telles que}
\note{sous $\{i\}$ une accolade porte $a'$, sous $\{i, j, k\}$ une accolade porte $A'$.}

\[
  a', b, c, \qquad b \vee c = A, \quad c \vee a' = B, \quad a' \vee b = C ;
\]
au-dessous, une flèche $\complement$ descend de $A$ vers $a$ (« pas OK ! »), de $B$ vers $b$ (« OK »), de $C$ vers $c$ (« OK »).

\page{114}
\note{p. 109 de l'auteur, chiffre récrit sur un 108.}
Après cette longue digression, il faudrait revenir à l'axiome des supports dans un atelier.

\textbf{Lemme 1.} Soit $\Sigma$ un spatios, $e_* = (e_i)_{i \in I}$ une famille d'élts de $\Sigma$ mutuellement disjoints, \struck{soit telle que} non nuls, et tels que
\[
  \bigvee_i e_i = 1_\Sigma .
\]
Conditions équivalentes
\begin{itemize}
  \item[(i)] L'ens.\ des $e_i$ est contenu dans un sous-ensemble \struck{$\mathfrak{B}$} de $\Sigma$ stable par Sup, Inf et $\complement$, \struck{$\vee$, $\wedge$, $\complement$, et qui} et qui, pour la structure induite \uncertain{et} $\complement$ induite, est une « alg.\ de Boole ».
  \item[(ii)] \struck{L'ensemble} Pour toute $J \in \mathfrak{P}(I)$, posant
  \[
    e_J = \operatorname*{Sup}_{i \in J} e_i ,
  \]
  l'application ainsi obtenue
  \[
    \mathfrak{P}(I) \longrightarrow \Sigma, \qquad J \longmapsto e_J
  \]
  (qui commute \uncertain{déjà} aux Sup quelconques) commute : $\complement$ (donc aussi aux Inf quelconques).
\end{itemize}
\marginal{NB cette application \uncertain{est} injective \ill{} $= \{i \in I \mid e_i \le e\}$}

\page{115}
\note{p. 110 de l'auteur.}
Sous ces conditions, \struck{\ill{}} l'image $\Sigma_e$ de cette application est le plus petit sous-ensemble $\Sigma'$ de $\Sigma$ satisfaisant aux conditions de (i).

\textbf{Cor.\ 1.} Mêmes données, mais on ne suppose pas que $e \overset{\text{déf}}{=} \operatorname*{Sup}_{i \in I} e_i$ soit égal à $1$, et on \uncertain{introduit} $e_0 = \complement e$. Pour que les conditions des lemmes soient satisfaites, \uncertain{il} \uncertain{suffit} \uncertain{que} pour
\[
  \bigl\{ \{e_i\}, e_0 \bigr\}
\]
(\uncertain{indexé} par $I \sqcup \{0\}$) il faut et il suffit \struck{qu'on ait}, pour toute partie $I'$ de $I$, posant $I'' = I \smallsetminus I'$, $e = e_{I'} \vee e_{I''}$ (\uncertain{de} \uncertain{sorte} que $e_0 = \complement e$), \struck{on ait} \struck{on} ait
\[
  \complement e_{I'} = e_{I''} \vee e_0
\]
(\uncertain{compatibilité} aux $\complement$), ce qui revient au même à
\[
  e_{I'} = e \wedge \complement e_{I''} .
\]

\textbf{Déf.} On dit alors que la \struck{\ill{}} \add{famille} $(e_i)_{i \in I}$ d'éléments mutuellement disjoints est « \textit{admissible} ».
\[
  \Sigma_{e_*} = \Sigma_{(e_i)} ,
\]
alg.\ de Boole engendrée, iso.\ à $\mathfrak{P}\bigl(I \sqcup \{e_0\}\bigr)$.

\page{116}
\note{p. 111 de l'auteur.}
Ainsi, le premier axiome sur les supports d'un atelier \struck{$A$} muni d'un magasin \struck{\ill{}} \add{associé} $\mathcal{F}$, sera \add{$\mathrm{Figat}(A)$} :

\textbf{\uncertain{Atmagsupp}~1.} Pour toute figure $F \in \mathcal{F}_c$, la famille des $(X^\circ)_{X \in F}$ dans $\Sigma_A$ est admissible.

Cet axiome n'est pas redondant, \ill{} \uncertain{on} \uncertain{le} \uncertain{voit} \uncertain{en} \uncertain{partant} d'un \uncertain{commun} \ill{} en prenant l'atelier \uncertain{associé} au spatios $\Sigma$, \struck{\ill{}} $\mathrm{Figelspat}_{\mathrm{fin}}(\Sigma)$ \struck{\ill{}} des figures élémentaires (spatiales) finies (\uncertain{disjointes}) de $\Sigma$, et le magasin \struck{des} \uncertain{associé} $\mathrm{Figspat}_{\mathrm{fin}}(\Sigma)$ des figures spatiales finies. La validité de l'axiome des supports revient alors, d'après le cor.\ 1, à ceci : si $e, e' \in \Sigma$ \struck{avec} $e \mathbin{|\circ|} e'$, \struck{et si $e'' = \complement(e \vee e')$,} \struck{alors}
\[
  \text{alors} \quad e' = (\complement e) \wedge (e \vee e') .
\]
\uncertain{Du} \uncertain{moins}, la validité des conditions du lemme 2 (p.~105) sur $\Sigma$.
\note{« p.~105 » renvoie à la pagination de l'auteur, c'est-à-dire à la p.~108 des archivistes (lemme 2 sur $b \le c$).}

\page{117}
\note{p. 112 de l'auteur.}
L'axiome des supports ne concerne directement que les données $\lhd$, $\mathbin{|\circ|}$ et $\mathcal{F} \subset \mathfrak{P}(A)$ sur $A$, et ne fait pas intervenir \uncertain{encore} $\overset{\circ}{\ll}$ ou $\ll$. Je \uncertain{conviens} de la \uncertain{convention} \uncertain{suivante}.

Considérons dans un spatios $\Sigma$ toutes les \struck{familles $(e_*, I) = (e_i)_{i \in I}$} \struck{$e_* = \{e_i\}_{i \in I}$} \add{parties} $A$, formées d'objets \struck{non} \uncertain{nuls} mutuellement disjoints. On peut les considérer comme des objets de $\mathrm{Figspat}(\Sigma)$, \struck{pour} la relation d'ordre \uncertain{étant} discrète (\uncertain{discrets}) \struck{sur} \struck{\ill{}} $A$. Nous allons écrire, \uncertain{pour}
\[
  A, B \in \mathrm{Figspatdisc}(\Sigma) ,
\]
\[
  B \ll A \iff \forall b \in B,\ \exists a \in A \text{ avec } b \le a
\]
(\textbf{NB} c'est aussi la relation \add{induite par} $\ll$ de $\mathrm{Figspat}$).

\struck{Cons} Pour un $a \in A$ donné, considérons l'ens.
\[
  B_a = \{ b \in B \mid b \le a \} ,
\]
formé des élts de $B$ (\uncertain{mut.}\ disj.) qui sont $\le a$. Soit $a' = \operatorname{Sup} B_a \le a$, et $a'' =$ \uncertain{$a \wedge \complement a'$}.

\page{118}
\note{p. 113 de l'auteur.}
\textbf{Lemme 2.} Avec les notations précédentes, les conditions suivantes sont équivalentes
\begin{itemize}
  \item[(i)] $A \cup B \subset \Sigma'$, où $\Sigma'$ est une partie de $\Sigma$ stable par Sup, Inf et $\complement$, et qui est une alg.\ de Boole pour l'ordre induit par $\Sigma$ et pour l'op.\ induite par $\complement$.
  \item[(ii)] On a
\end{itemize}
\begin{itemize}
  \item[a)] $\forall a \in A$, on a $a = a' \vee a''$ ;
  \item[b)] \struck{La famille formée des $b \in B$ et parties de $B$} L'ensemble
  \[
    \widehat{B} = B \cup \{ a'' \mid a \in A \text{ t.q. } a'' \ne 0 \}
  \]
  est admissible.
\end{itemize}
Sous ces conditions, l'alg.\ de Boole $\Sigma_{\widehat{B}}$ engendrée par $\widehat{B}$ (cf.\ déf.\ p.~110), \struck{est aussi l'} iso.\ à $\mathfrak{P}(\widehat{B})$, est aussi le plus petit des sous-ens.\ $\Sigma'$ satisfaisant (i).
\note{« déf.\ p.~110 » renvoie à la pagination de l'auteur : la définition des familles admissibles, p.~115 des archivistes.}

\textbf{Déf.} On dit donc que $\{A, B\}$ est \textit{admissible}.

\textbf{Axiome \uncertain{Atmagsupp}~2.} Pour \struck{$F$} \add{toute} $F \ll G$ avec $F, G \in \mathcal{F}$, posant
\[
  B = \{ X^\circ \mid X \in G \}, \qquad A = \{ X^\circ \mid X \in F \} ,
\]
\note{ainsi sur la page : $B$ est indexé par $G$ et $A$ par $F$, alors que $F \ll G$ et que $B \ll A$ dans la notation de la p.~117.}

\page{119}
\note{p. 114 de l'auteur.}
\struck{Les deux familles parties} $\{A, B\}$ est une partie admissible de parties de $\Sigma_A$.

Ceci \uncertain{n'est} pas \uncertain{assez} : \uncertain{il} \uncertain{faudrait} \uncertain{pouvoir} traiter le cas d'un « drapeau » quelconque dans $(\mathcal{F}, \ll)$
\[
  F_0 \ll F_1 \ll \dots \ll F_n
\]
donnant naissance dans $\Sigma = \Sigma_A$ à un drapeau similaire dans $\mathrm{Figspat}(\Sigma)$, mais qui \struck{\ill{}} \uncertain{n'a} \uncertain{plus} \uncertain{d'}interprétation, \uncertain{par} \uncertain{oubli} des structures d'ordre supplémentaires sur \struck{\ill{}} chacun des $F_i$, comme un drapeau dans $\mathrm{Figspatdisc}(\Sigma)$
\[
  A_0 \ll A_1 \ll \dots \ll A_n .
\]

Pour exprimer l'axiome pertinent, \struck{disons de façon générale} \uncertain{généralisons} la

\textbf{Définition.} Soit $\Sigma$ un spatios, \struck{une} \struck{et} $A$ une partie de $\Sigma$. Nous dirons que $A$ est une partie \uncertain{commensurable} de $\Sigma$, \uncertain{ou} que \struck{\ill{}}

\page{120}
\note{p. 115 de l'auteur.}
les $a \in A$ sont « \uncertain{commensurables} dans leur ensemble » (\textbf{NB} c'est plus fort que d'exiger qu'ils le soient $2 : 2$), si $A \subset \Sigma'$, $\Sigma'$ \struck{un} un « sous-spatios » (stable par Sup, $\complement$ \uncertain{et} par Inf) qui est une alg.\ de Boole.

Ceci posé, étant donné un drapeau
\[
  A_0 \ll A_1 \ll \dots \ll A_n ,
\]
nous dirons qu'il est commensurable, si
\[
  A = \bigcup_{1 \le i \le n} A_i
\]
l'est.

\struck{On pourrait \ill{} l'exp}

\textbf{Lemme 3.} Pour tout $1 \le i \le n$, et $a_i \in A_i$, considérons $a''_i$ défini comme dans lemme 2, pour $A_{i-1} \ll A_i$. Soit
\[
  A''_i = \{ a''_i \mid a_i \in A_i \text{ t.q. } a''_i \ne 0_\Sigma \} .
\]
\struck{Soit} On pose aussi
\[
  A''_0 = A_0 , \qquad
  A = \bigcup_{0 \le i \le n} A''_i .
\]
Alors les $A''_i$ sont formés d'élts \add{non nuls} mut.\ disj., et ils sont mutuellement disjoints entre eux,

\note{l'énoncé du lemme 3 se poursuit au-delà de ce lot.}

\end{document}
